% Source: source/普通高考/2020/2020全国1文(河南,河北,山西,江西,湖北,湖南,广东,安徽,福建).pdf, page 1.
% The source flowchart is vector line art; pdfimages found no matching embedded flowchart.
% Original comparison crop: tmp/redraw_flowchart_2020_national_paper_1_liberal/q9_orig_source.png,
% cropped from a no-grid page render after grid100/grid50 localization.
% Elements: start/end rounded rectangles, input/output parallelograms, two process rectangles,
% decision diamond S<=100, directed connectors, and the 是/否 branch labels.
% Node→directed-edge list: 开始→输入(n=1,S=0)→S=S+n→S<=100;
% S<=100 的 是→n=n+2→S=S+n；S<=100 的 否→输出 n→结束.
% solid segments: all node borders and connectors; dashed segments: none.
% labels: every node is locally zero-indented and horizontally/vertically centered;
% 是 marks the left loop branch and 否 marks the downward output branch.
\begin{tikzpicture}[>=stealth, line width=0.45pt]
  % The source is a schematic flowchart; coordinates preserve its topology and labels.
  \tikzset{flow/.style={draw=black,line width=0.45pt,align=center,inner sep=0.6pt,
    execute at begin node={\setlength{\parindent}{0pt}}}}
  \tikzset{every node/.style={font=\normalsize,align=center,inner sep=0pt,
    execute at begin node={\setlength{\parindent}{0pt}}}}
  \node[flow, rounded corners=2pt, minimum width=1.05cm, minimum height=0.62cm]
    (start) at (0,7.55) {开始};
  % Input parallelogram is drawn around a centered text node to keep the source margins.
  \draw[flow] (-1.80,6.90) -- (1.90,6.90) -- (1.70,6.20) -- (-2.00,6.20) -- (-1.80,6.90);
  \node (input) at (0,6.55) {输入 \(n=1,S=0\)};

  \node[flow, minimum width=1.85cm, minimum height=0.58cm] (inc)
        at (-2.55,5.32) {\(n=n+2\)};
  \node[flow, minimum width=1.95cm, minimum height=0.58cm] (sum)
        at (0,5.32) {\(S=S+n\)};

  \node[flow, diamond, aspect=2.8, minimum width=3.25cm, minimum height=0.78cm]
        (test) at (0,4.10) {\(S\leq 100\)};
  \node[flow, trapezium, trapezium left angle=72, trapezium right angle=108,
        minimum width=1.50cm, minimum height=0.60cm] (output) at (0,2.85)
        {输出 \(n\)};
  \node[flow, rounded corners=2pt, minimum width=1.05cm, minimum height=0.62cm]
        (finish) at (0,1.65) {结束};

  % Downstream path and loop connectors.
  \draw[->,flow] (start.south) -- (0,6.90);
  \draw[->,flow] (0,6.20) -- (sum.north);
  \draw[->,flow] (sum.south) -- (test.north);
  \draw[->,flow] (test.south) -- (output.north)
        node[pos=0.50, right=2pt] {否};
  \draw[->,flow] (output.south) -- (finish.north);

  % Yes branch: diamond left tip -> increment -> return to the sum step.
  \draw[flow] (test.west) -- (-2.55,4.10) -- (-2.55,4.70)
        node[pos=0.02, above right=1pt] {是};
  \draw[->,flow] (-2.55,4.70) -- (inc.south);
  % Route the return above both process boxes, preserving clearances.
  \draw[flow] (inc.north) -- (-2.55,5.90);
  % Stop the return arrow before the shared main-shaft junction; finish with a plain shaft.
  \draw[->,flow] (-2.55,5.90) -- (-0.45,5.90);
  \draw[flow] (-0.45,5.90) -- (0,5.90) -- (sum.north);
\end{tikzpicture}
